\section{Introduction}\label{sec:intro}

Catastrophic health expenditure---defined as out-of-pocket (OOP) spending exceeding a fixed share of household resources---remains a central concern for health systems in developing countries \citep{Xu2003, Wagstaff2003}. In Peru, the coexistence of a subsidized public insurance scheme for the poor (Seguro Integral de Salud, SIS), a contributory scheme for formal workers (EsSalud), and a large uninsured population creates a fragmented landscape in which financial protection against health shocks varies sharply across the income distribution \citep{Petrera2013, Francke2013}. Despite two decades of coverage expansion under the Universal Health Coverage framework established by Law 29344 (Ley Marco de Aseguramiento Universal en Salud, 2009), evidence on how effectively these schemes reduce financial risk across different segments of the spending distribution remains limited.

The standard approach in the CHE literature estimates the prevalence of households exceeding a binary threshold---typically 10\%, 25\%, or 40\% of total consumption \citep{Xu2003}---and examines its correlates through logistic regression at the conditional mean. This approach obscures two dimensions of heterogeneity. First, it collapses the distribution of OOP burden into a single binary indicator, discarding information about spending intensity. Second, mean-based regression assumes that covariates shift the spending distribution uniformly, a restriction that quantile regression relaxes \citep{Koenker1978, Koenker2005}. Moreover, the large fraction of households reporting zero OOP spending---a pervasive feature of health expenditure data in developing countries---renders standard quantile regression degenerate at lower quantiles without appropriate adjustments \citep{Powell1986, Manning1987}.

This paper applies conditional quantile regression to Peru's 2024 National Household Survey (Encuesta Nacional de Hogares, ENAHO), interacting insurance status with consumption quintiles to examine whether the association between insurance coverage and OOP health expenditure share varies across the conditional spending distribution. We estimate models at $\tau = \{0.50, 0.75, 0.90\}$---quantiles at which the conditional quantile function is positive and informative given the 44.1\% zero-OOP mass in our sample---and complement the analysis with a two-part model that separately identifies the extensive margin (probability of any OOP spending) and the intensive margin (spending conditional on positive expenditure).

Our analysis proceeds in four steps. First, we document the distributional features of OOP health expenditure in Peru, highlighting the zero-spending mass and the pronounced right skew of positive spending. Second, we estimate survey-weighted OLS and conditional quantile regressions with insurance--quintile interactions, demographic controls, and region fixed effects, using cluster-bootstrapped standard errors at the primary sampling unit (PSU) level. Third, we decompose the total insurance association by quintile to identify where in the distribution insurance coverage is most strongly associated with reduced OOP burden. Fourth, we subject our findings to 25 robustness checks spanning alternative samples, CHE thresholds, consumption floor sensitivity, control sensitivity, and placebo tests.

Three provisional findings emerge, subject to the data and inference limitations detailed below. First, the base SIS coefficient is statistically indistinguishable from zero across all quantiles when evaluated at the lowest consumption quintile (Q1), consistent with the interpretation that SIS enrollees in Q1 either access free care or forgo care altogether. Second, the SIS--quintile interaction terms are negative and increasingly significant at higher quantiles and higher quintiles: at $\tau = 0.90$, SIS membership in Q5 is associated with a reduction of 1.94 percentage points in OOP share ($p < 0.01$ with 200-replication bootstrap; significance is provisional pending 999-replication confirmation). This monotonic gradient---where the insurance association \textit{increases} with income conditional on enrollment---is more accurately characterized as a regressive protection pattern than as a middle-income squeeze. Third, consumption quintile gradients steepen monotonically across quantiles, with the Q5 coefficient rising from near zero at $\tau = 0.10$ to 5.16 percentage points at $\tau = 0.90$. However, this gradient is partly mechanical: consumption quintiles are constructed from \texttt{GASHOG2D}, which includes health expenditure, creating a positive correlation between quintile rank and OOP share. Reconstruction of quintiles from non-health consumption is required to isolate the behavioral gradient from the mechanical component.

We restrict the primary analysis to the SIS--uninsured comparison. Our EsSalud sample ($N = 648$ household heads, approximately 1.9\% of the sample) is implausibly small relative to national EsSalud coverage rates of approximately 30\%, likely reflecting a variable construction error in the household-head-based insurance assignment that discards non-head EsSalud members. This misclassification contaminates the uninsured reference group with households that have EsSalud coverage through non-head members. EsSalud results are reported in supplementary tables but are excluded from substantive interpretation.

We interpret the monotonic SIS--quintile gradient cautiously. SIS is means-tested through the SISFOH (Sistema de Focalizaci\'on de Hogares), and its presence in upper quintiles reflects well-documented targeting errors and leakage \citep{Cortez2008, Francke2013}. The larger SIS--quintile interactions at Q4--Q5 may reflect a composition effect---non-poor SIS enrollees who access care more effectively than their poor counterparts---rather than evidence that SIS provides differential financial protection by design. Our descriptive framework cannot distinguish between these mechanisms.

These findings contribute to three literatures. First, we add to the growing body of work applying quantile regression to health expenditure analysis \citep{Jones2013, Machado2019}, demonstrating that the conditional correlates of OOP burden are quantile-dependent in a middle-income country context. Second, we contribute to the literature on Peru's health financing system \citep{EspinozaPajuelo2024, Petrera2013}, providing the first distributional analysis of CHE determinants using the most recent ENAHO wave. Third, we inform the policy debate on universal coverage design by showing that average CHE rates mask substantial heterogeneity in financial risk across both the income and spending distributions.

We emphasize throughout that our estimates represent conditional associations, not causal effects. Insurance enrollment in Peru is non-random: SIS targets the poor through means testing, while EsSalud covers formal-sector workers \citep{Wagstaff2010, Acharya2012}. We note further that the 2024 ENAHO is collected during the post-pandemic recovery period; Peru experienced among the world's highest COVID-19 excess mortality rates, with substantial disruption to health-seeking behavior. Health expenditure patterns in 2024 may reflect post-pandemic adjustment rather than steady-state behavior, a caveat that limits the generalizability of our findings.

The remainder of the paper is organized as follows. Section~\ref{sec:literature} reviews the relevant literature. Section~\ref{sec:data} describes the data and variable construction. Section~\ref{sec:strategy} presents the empirical strategy. Section~\ref{sec:results} reports the main results. Section~\ref{sec:robustness} discusses robustness checks. Section~\ref{sec:conclusion} concludes.
